% !TeX program = xelatex
\documentclass[11pt,oneside]{article}
\usepackage{ProblemsetStyle}

\hypersetup{
    pdftitle={20260826 Problem set - Solutions},
    pdfauthor={김정우},
    pdfsubject={Science Quiz mathematics practice problem set solutions},
    pdfcreator={XeLaTeX}
}

\begin{document}

\problemsettitle{20260826 Problem set}

\begin{problemsolution}{1}
수열
\[
    b_n=2\cos\left(\frac{n\pi}{4}\right)
\]
을 생각하자. 이 수열은 $a_n$과 같은 초기값을 가지며,
\[
    b_{n+2}
    =
    2\cos\left(\frac{\pi}{4}\right)b_{n+1}-b_n
    =
    \sqrt2\,b_{n+1}-b_n
\]
을 만족한다. 따라서 $a_n=b_n$이다. $2026\equiv2\pmod8$이므로
\[
    \boxed{
        a_{2026}
        =
        2\cos\left(\frac{2026\pi}{4}\right)
        =0
    }.
\]
\end{problemsolution}

\begin{problemsolution}{2}
배각공식
\[
    (2\cos\theta)^2-2=2\cos(2\theta)
\]
에 의해 귀납적으로
\[
    a_n
    =
    2\cos\left(\frac{2^{n+1}\pi}{7}\right)
\]
을 얻는다. $2$의 거듭제곱을 $14$로 나눈 나머지는 주기 $3$을 가지며,
\[
    2^{2027}\equiv4\pmod{14}.
\]
따라서
\[
    \boxed{a_{2026}=2\cos\left(\frac{4\pi}{7}\right)}.
\]
\end{problemsolution}

\begin{problemsolution}{3}
모든 $n$에 대하여 $f_n(1)=1$이다. 로그를 취하면
\[
    \ln f_{n+1}(x)=f_n(x)\ln x.
\]
미분하여 $x=1$을 대입하면
\[
    f_{n+1}'(1)=f_n(1)=1
\]
이므로, 귀납적으로 $f_n'(1)=1$이다.

이제
\[
    g(x)=\ln f_{n+1}(x)=f_n(x)\ln x
\]
라 두자. 그러면
\[
    g'(1)=1,
\]
그리고
\[
    g''(1)
    =
    2f_n'(1)-f_n(1)
    =1.
\]
$f_{n+1}=e^g$이므로
\[
    f_{n+1}''(1)
    =
    f_{n+1}(1)\bigl(g'(1)^2+g''(1)\bigr)
    =2.
\]
따라서
\[
    \boxed{f_{2026}''(1)=2}.
\]
\end{problemsolution}

\begin{problemsolution}{4}
Fibonacci sequence를 $F_0=0$, $F_1=1$로 두면 그 ordinary generating function은
\[
    \frac{1}{1-x-x^2}
    =
    \sum_{n=0}^{\infty}F_{n+1}x^n
\]
이다. 따라서 $x^{2026}$의 coefficient는 $F_{2027}$이고,
\[
    \boxed{f^{(2026)}(0)=2026!\,F_{2027}}.
\]
\end{problemsolution}

\begin{problemsolution}{5}
\[
    e^x\sin x
    =
    \operatorname{Im}\bigl(e^{(1+i)x}\bigr)
\]
이므로 요구된 값은
\[
    \operatorname{Im}\bigl((1+i)^{2026}\bigr)
\]
이다. 한편
\[
    (1+i)^{2026}
    =
    2^{1013}e^{i2026\pi/4}
    =
    2^{1013}i.
\]
따라서
\[
    \boxed{2^{1013}}.
\]
\end{problemsolution}

\begin{problemsolution}{6}
2026 International Congress of Mathematicians의 개최지는
\[
    \boxed{\text{Philadelphia, United States}}
\]
이다. 네 명의 Fields Medal 수상자는
\[
    \boxed{\text{Yu Deng, John Pardon, Jacob Tsimerman, Hong Wang}}
\]
이다.
\end{problemsolution}

\begin{problemsolution}{7}
해당 수학자는
\[
    \boxed{\text{Hong Wang}}
\]
이다. Hong Wang은
\[
    \boxed{\text{Fields Medal을 받은 세 번째 여성}}
\]
이 되었다.
\end{problemsolution}

\begin{problemsolution}{8}
2026 Abel Prize 수상자는
\[
    \boxed{\text{Gerd Faltings}}
\]
이다. Faltings는 1986년 Fields Medal 수상자이기도 하다.
\end{problemsolution}

\end{document}
